% 第5章图：消去x2时的函数更新与图更新（四变量例子，作用域{1,2,3},{2,3,4}）。
\begin{tikzpicture}[x=1mm,y=1mm,>=Stealth,
  v/.style={circle,draw=FigureStroke,fill=FigureBlueFill,line width=.8pt,minimum size=6.5mm,inner sep=0pt},
  cur/.style={v,fill=FigureYellowFill,draw=FigureYellow},
  ghost/.style={circle,draw=FigureGrid,dashed,line width=.6pt,minimum size=6.5mm,inner sep=0pt,text=FigureGrid},
  e/.style={FigureStroke,line width=.8pt},
  fbox/.style={rectangle,rounded corners=1.8pt,draw=FigureStroke,fill=FigureGreenFill,line width=.7pt,inner sep=1.5pt},
  note/.style={font=\FigureNoteFont\fontsize{7.5}{9}\selectfont,text=FigureMuted}]
\begin{scope}
 \node[v] (l1) at (0,10) {$1$};
 \node[cur] (l2) at (14,22) {$2$};
 \node[v] (l3) at (14,0) {$3$};
 \node[v] (l4) at (28,10) {$4$};
 \draw[e] (l1)--(l2) (l1)--(l3) (l2)--(l3) (l2)--(l4) (l3)--(l4);
 \node[fbox] at (14,-11) {$\phi_A(x_1,x_2,x_3)$\quad$\phi_B(x_2,x_3,x_4)$};
 \node[note] at (14,-18) {当前函数集合};
\end{scope}
\draw[->,FigureStroke,line width=.9pt] (39,10)--(55,10) node[midway,above,font=\figurefont\footnotesize]{消去$x_2$};
\node[note,align=center] at (47,4) {合并含$2$的函数\\对$x_2$取最小};
\begin{scope}[xshift=66mm]
 \node[v] (r1) at (0,10) {$1$};
 \node[ghost] (r2) at (14,22) {$2$};
 \node[v] (r3) at (14,0) {$3$};
 \node[v] (r4) at (28,10) {$4$};
 \draw[e] (r1)--(r3) (r3)--(r4);
 \draw[FigureGreen,line width=1pt,dashed] (r1)--(r4);
 \node[font=\FigureNoteFont\fontsize{7.5}{9}\selectfont,text=FigureGreen,above] at (14,10.5) {填充边};
 \node[fbox] at (14,-11) {$\psi_2(x_1,x_3,x_4)$};
 \node[note] at (14,-18) {消元后的函数集合};
\end{scope}
\end{tikzpicture}
